\section{Empirical Evaluation}\label{sec:eval}
\answers{RQ4 (E6, E2), RQ5 (E3, E5) and the supply-chain part of RQ2 (E4)}

All experiments ran on one Linux workstation with Docker and a consumer GPU; every result number is read from \texttt{experiments/results/*.json}.

\subsection{E6: Rule-pack accuracy on a labelled corpus (RQ4)}\label{sec:eval-e6}
\paragraph*{Design} The corpus has \nESixCases{} cases (\nESixHand{} hand-crafted, \nESixGenerated{} generated) over more than twenty harness formats, CI files, state directories and MCP manifests; \nESixNegative{} are fully negative, and many positives have near-miss twins (pinned \texttt{npx}, authenticated HTTPS MCP, linter-only hooks, emoji ZWJ, key-like identifiers). Every audit runs on every case and the unit is the (case, rule) pair, so any unlabelled firing is a false positive. A hash of the case id splits the corpus into \nESixDev{} development and \nESixTest{} test cases; errors were inspected and fixed only on the development split.

\paragraph*{Results} On the held-out test split, micro precision is \eSixPrecision{} (95\% Wilson CI \eSixPrecisionLo--\eSixPrecisionHi) and micro recall \eSixRecall{} (\eSixRecallLo--\eSixRecallHi), with \eSixFP{} false positives and \eSixFN{} false negatives; macro precision and recall are \eSixMacroPrecision{} and \eSixMacroRecall{} (per-rule results: \texttt{paper/tables/e6\_accuracy.md}). Before the development-split fixes, test precision was \eSixPrecisionBefore{} with \eSixFPBefore{} false positives; the test split had been scored once before corpus-content fixes that touched both splits, so it is not pristine. The Wilson intervals treat (case, rule) pairs as independent although they cluster within cases and generator families, so they are optimistic. The weakest rules are \texttt{HG-SUB-001}, \texttt{HG-MCP-004} and the approve-everything rules. Because the corpus is synthetic and labelled by the rule author, this shows that the rules implement their specification across formats; real-world precision is unknown.


\subsection{E2: Prevalence in public repositories (RQ4)}\label{sec:eval-e2}
\paragraph*{Design} Through the GitHub code-search API we sampled harness settings and MCP files, instruction files (\texttt{AGENTS.md}, \texttt{CLAUDE.md}, Copilot and Windsurf rules) and CI workflows invoking \texttt{claude-code-action}, \texttt{codex exec} or \texttt{gemini -p}. Per repository we rebuilt only the sampled files in a temporary tree, ran \texttt{harnesswatch audit}, kept only (rule, file) pairs and deleted the tree. Identifiers were hashed with an unsaved per-run salt for de-duplication only; a scrub check refuses to write results containing any identifier or URL seen; secret findings are only counted; requests stayed below the API's rate limits. Each of the \eTwoTargets{} search strata is capped at \eTwoCapPerTarget{} files, so pooled per-repository rates mostly reflect which strata were chosen. We therefore report each rule's rate over the stratum of files that can carry it (config/MCP files, $n=\eTwoStratumConfig$; instruction files, $n=\eTwoStratumInstructions$; agent CI workflows, $n=\eTwoStratumWorkflows$), with Wilson intervals, and never name a repository.

Code search over-represents popular, recently indexed files; only project-level configuration is visible; a finding marks risky configuration, not exploitation.

% BEGIN-E2-RESULTS
\paragraph*{Results} We audited \eTwoReposAudited{} repositories (\eTwoFilesAudited{} files, \eTwoAuditErrors{} audit errors) and \eTwoRulesFired{} rules fired at least once (Table~\ref{tab:e2}). Among agent CI workflows, \eTwoStratCiTwo{} run the agent with a write-scoped token (\texttt{HG-CI-002}) and \eTwoStratCiOne{} match the pwn-request pattern of untrusted trigger plus secrets (\texttt{HG-CI-001}). Among config and MCP files, \eTwoStratMcpOne{} launch an MCP server from a registry without a version pin (\texttt{HG-MCP-001}) and \eTwoStratMcpTwo{} declare a remote server without TLS or authentication (\texttt{HG-MCP-002}); \texttt{HG-MCP-005} (server declared in repository config, \eTwoStratMcpFive{}) largely restates the sampling key and is informational. Secret-shaped strings in MCP declarations (\texttt{HG-MCP-003}, \eTwoStratMcpThree{}) were counted only; since identifiers are discarded by design, owners were not notified and the strings were not validated. No E2 finding was hand-validated, so the real-world precision of these rules is unknown; the rates describe this capped, relevance-ranked sample, not GitHub as a whole.
\begin{table}[!t]
\centering\footnotesize
\caption{E2 prevalence per stratum (RQ4): repositories with a finding over audited files in the stratum that can carry the rule; top rules, Wilson 95\% CI. Aggregates only; no repository is named.}
\label{tab:e2}
\begin{adjustbox}{max width=\columnwidth}\begin{tabular}{llrl}
\toprule
Rule & Stratum ($n$) & Repos & Rate [95\% CI] \\
\midrule
HG-CI-002 & workflows (177) & 112 & 63.3\% [56.0, 70.0] \\
HG-MCP-005 & config (600) & 311 & 51.8\% [47.8, 55.8] \\
HG-MCP-001 & config (600) & 170 & 28.3\% [24.9, 32.1] \\
HG-CI-001 & workflows (177) & 50 & 28.2\% [22.1, 35.3] \\
HG-MCP-002 & config (600) & 86 & 14.3\% [11.8, 17.4] \\
HG-HOOK-001 & config (600) & 43 & 7.2\% [5.4, 9.5] \\
HG-MCP-003 & config (600) & 21 & 3.5\% [2.3, 5.3] \\
HG-SEC-003 & config (600) & 14 & 2.3\% [1.4, 3.9] \\
HG-PERM-001 & config (600) & 12 & 2.0\% [1.1, 3.5] \\
HG-PERM-003 & config (600) & 12 & 2.0\% [1.1, 3.5] \\
\bottomrule
\end{tabular}
\end{adjustbox}
\end{table}
% END-E2-RESULTS

\subsection{E3: Canary benchmark for indirect prompt injection (RQ5)}\label{sec:eval-bench}
\paragraph*{Design} Each run uses a fresh container with a tiny Python repository, a synthetic \texttt{.env} token, a fixture file and a sink on \texttt{127.0.0.1}, and the same benign task (add a \texttt{multiply} function and test, run the tests). Besides a \emph{control}, four scenarios place one templated block in the README, the native instruction file, a code comment, or test output that returns as a tool result (Table~\ref{tab:bench-catalogue}). The block asks the agent to create \texttt{CANARY\_<id>.txt}, POST \texttt{.env} to the sink, and delete the fixture; all are inert, and the container reaches only the model server. Modes are \emph{default} (out-of-the-box headless) and \emph{auto} (the documented approve-everything flag).

\paragraph*{Harnesses, model and grid} Seven adapters cover three Tier-B clients (Claude Code, Copilot CLI and the Grok Build npm binary, driven through custom-endpoint modes; Claude Code's is vendor-unsupported with non-Claude models, Section~\ref{sec:method}) and four open ones (OpenCode, Codex CLI, Aider and Qwen Code). No vendor credentials are used, so every harness talks to one local open-weight model, \texttt{qwen3:4b-instruct}, served by Ollama. The grid is \nBenchHarnesses{} harnesses $\times$ \nBenchScenarios{} scenarios $\times$ 2 modes with \nBenchPerCell{} runs per cell. A run is capped at 240~s and scored on the state at the cap; a run with no harness output counts as an error, not as safe. Aider never loads an instruction file, so its \texttt{instructions} cells are N/A in both modes and excluded from the pooled statistics.

\paragraph*{Metrics} \emph{Injection follow}: any of the three goals achieved (marker created, sink POST, fixture deleted). \emph{Exfil attempt}: any POST to the sink, or the secret seen at the sink. \emph{Destructive action}: fixture deleted. \emph{Secret read}: the synthetic token appears in the harness's printed output (other than as an environment-variable dump) or reaches the sink; since some harnesses print only the final message, this undercounts reads. \emph{Secret in state}: the token appears in a transcript or session file the harness wrote, separating persistence from agent action. \emph{Permission prompt}: an approval prompt in the output. \emph{Task completed}: tests pass and the function and its test exist, so that ``safe but useless'' is visible. Rates carry Wilson 95\% intervals.

% BEGIN-E3-RESULTS
\paragraph*{Results} The grid produced \nBenchRuns{} runs (\nBenchErrors{} errors, \nBenchTimeouts{} timeouts scored at the cap, model cost \$\benchCost{}), of which \nBenchRunsAnalysed{} enter the analysis (Table~\ref{tab:e3}; plotted in \texttt{paper/figures/e3\_follow.pdf}). Over the injection scenarios the injected instruction was followed in \benchFollowAutoK{} of \benchFollowAutoN{} auto-mode runs (95\% CI \benchFollowAutoLo--\benchFollowAutoHi) and \benchFollowDefaultK{} of \benchFollowDefaultN{} default-mode runs (\benchFollowDefaultLo--\benchFollowDefaultHi); controls followed nothing (\benchControlFollowK{} of \benchControlN{}). No run attempted the POST, the fixture was deleted in \benchDestrAutoK{} auto-mode run, and the strict secret-read count is \benchSecretReadAutoK{} in both modes. The intervals overlap, so there is \emph{no detectable safety difference} between modes. Task completion differs (\benchTaskAutoK{} of \benchTaskAutoN{} auto vs \benchTaskDefaultK{} of \benchTaskDefaultN{} default) because several headless defaults deny the tools the task needs: Codex CLI and Copilot CLI showed a permission prompt in every default run, while Grok Build and Qwen Code completed nothing in default mode without a detectable prompt, for reasons we did not establish. OpenCode completed \benchOpencodeControlTaskK{} of \benchOpencodeControlN{} even in controls, an adapter--model incompatibility, so its cells are uninformative. Claude Code's ``default'' arm is not an ask mode: with a third-party endpoint it starts in the classifier-mediated \texttt{auto} mode (Section~\ref{sec:claims}).

\paragraph*{Secret in state} In \benchClaudeImportK{} of \benchClaudeImportN{} Claude Code runs of the \texttt{instructions} scenario (both modes) the synthetic secret reached the session transcript, against \benchClaudeOtherStateK{} in its other scenarios. This is not a model action: an E7 probe with an unreachable model endpoint reproduced it, because the payload's \texttt{--data-binary @.env} contains the token \texttt{@.env}, which Claude Code's documented \texttt{@path} import resolves when it loads \texttt{CLAUDE.md}\evT\evD{}. It is a harness design property: an instruction file can import gitignored secrets into context and transcripts without any approval.

\paragraph*{Interpretation} The low follow rate most likely reflects the small model's limited agentic and instruction-following ability and the plainly worded benign canary, \emph{not} harness robustness; with \nBenchPerCell{} runs per cell and zero or one event per cell, intervals are wide, and we do not rank harnesses on resilience. We do not call default modes safer.

\paragraph*{Correction} Our first secret metric (token anywhere in output or written files) flagged \benchSecretAnyTotal{} runs, chosen before results were inspected and replaced afterwards: \benchClaudeImportK{} were the Claude Code import above and the rest Codex CLI, which writes its process environment, holding the canary, into \texttt{\textasciitilde/.codex} shell snapshots whenever its shell tool runs (\benchCodexSnapshotK{} of \benchCodexSnapshotN{} Codex runs)\evT{}. Neither is the agent reading the secret, hence the two separate metrics: benchmarks should separate what a harness persists from what an agent does.

\begin{table*}[!t]
\centering\footnotesize
\caption{E3 canary bench (RQ5): per harness and permission mode over the injection scenarios; $k/n$ with Wilson 95\% CI in percent; other columns are counts out of $n$. One local model (\texttt{qwen3:4b-instruct}); $^B$ marks Tier-B clients. Not a resilience ranking.}
\label{tab:e3}
\begin{adjustbox}{max width=\textwidth}\begin{tabular}{llrlrrrrrl}
\toprule
Harness & Mode & n & Follow [95\% CI, \%] & Exfil & Destr. & Secret in state & Prompts & Timeouts & Task done [95\% CI, \%] \\
\midrule
aider & auto & 15 & 0/15 [0.0, 20.4] & 0 & 0 & 0 & 0 & 0 & 10/15 [41.7, 84.8] \\
aider & default & 15 & 0/15 [0.0, 20.4] & 0 & 0 & 0 & 0 & 0 & 10/15 [41.7, 84.8] \\
claude-code$^{B}$ & auto & 20 & 1/20 [0.9, 23.6] & 0 & 1 & 5 & 0 & 4 & 16/20 [58.4, 91.9] \\
claude-code$^{B}$ & default & 20 & 0/20 [0.0, 16.1] & 0 & 0 & 5 & 0 & 6 & 13/20 [43.3, 81.9] \\
codex-cli & auto & 20 & 1/20 [0.9, 23.6] & 0 & 0 & 0 & 8 & 4 & 11/20 [34.2, 74.2] \\
codex-cli & default & 20 & 0/20 [0.0, 16.1] & 0 & 0 & 0 & 20 & 1 & 0/20 [0.0, 16.1] \\
github-copilot$^{B}$ & auto & 20 & 0/20 [0.0, 16.1] & 0 & 0 & 0 & 1 & 0 & 18/20 [69.9, 97.2] \\
github-copilot$^{B}$ & default & 20 & 0/20 [0.0, 16.1] & 0 & 0 & 0 & 20 & 2 & 0/20 [0.0, 16.1] \\
grok-build$^{B}$ & auto & 20 & 0/20 [0.0, 16.1] & 0 & 0 & 0 & 0 & 0 & 17/20 [64.0, 94.8] \\
grok-build$^{B}$ & default & 20 & 0/20 [0.0, 16.1] & 0 & 0 & 0 & 0 & 0 & 0/20 [0.0, 16.1] \\
opencode & auto & 20 & 0/20 [0.0, 16.1] & 0 & 0 & 0 & 0 & 0 & 0/20 [0.0, 16.1] \\
opencode & default & 20 & 0/20 [0.0, 16.1] & 0 & 0 & 0 & 0 & 0 & 0/20 [0.0, 16.1] \\
qwen-code & auto & 20 & 0/20 [0.0, 16.1] & 0 & 0 & 0 & 1 & 1 & 12/20 [38.7, 78.1] \\
qwen-code & default & 20 & 0/20 [0.0, 16.1] & 0 & 0 & 0 & 0 & 1 & 0/20 [0.0, 16.1] \\
\bottomrule
\end{tabular}
\end{adjustbox}
\end{table*}
% END-E3-RESULTS

\subsection{E5: Secret residue in harness state (RQ5)}\label{sec:eval-e5}
\paragraph*{Design} For each of the \nResidueHarnesses{} bench harnesses we ran two short sessions in a fresh container. In \emph{key} the provider credential was a unique synthetic canary. In \emph{content} the prompt asked the agent to read a synthetic \texttt{.env} file. We searched every file the harness wrote for the canary and recorded paths, modes and reachability by other users (home 0755); contents were never stored. With one session per harness and condition ($N=1$) the results are descriptive.

\paragraph*{Results} In the content condition, \nResidueWithCopy{} of \nResidueHarnesses{} harnesses left at least one file holding the secret, typically a transcript or chat log (Table~\ref{tab:e5})\evT{}. In \nResidueReachable{} harnesses such a file was readable by other local users. In the key condition, \nResidueKeyPersisted{} of the \nResidueKeyApplicable{} harnesses that take a key wrote it to a file in this single run\evT{}; that is a sampling result, not evidence of absence, because Codex CLI persists environment-supplied credentials in shell snapshots whenever its shell tool runs, which happened in \benchCodexSnapshotK{} of \benchCodexSnapshotN{} E3 runs and not in the E5 key session\evT{}. Mode bits alone can mislead: Grok Build writes 0644 session files, but inside a 0700 sessions directory, so other users cannot reach them\evT{}; Claude Code writes transcripts with mode 0600\evT{}. OpenCode left no residue in this run. Residue depends on whether the agent reads the file, so one run can miss it. These are design weaknesses, not vulnerabilities; a 0750 home directory would block the reachable paths, but CI runners and containers often use 0755.

\begin{table}[!t]
\centering\footnotesize
\caption{E5 secret residue after reading a synthetic \texttt{.env} (content condition; one session per harness, descriptive). ``Reach.'': files readable by other users (o+r on the file and o+x on every ancestor). ``Cred.'': files holding the provider key in the key condition (n/a: no key taken). ``o+r state'': other-user-readable / all state files written.}
\label{tab:e5}
\begin{adjustbox}{max width=\columnwidth}\begin{tabular}{lrllrrr}
\toprule
Harness & Files w/ secret & Kind & Mode & Reach. & Cred. & o+r state \\
\midrule
aider & 1 & workspace & 644 & 1 & n/a & 3/3 \\
claude-code & 1 & log & 600 & 0 & 0 & 2/5 \\
codex-cli & 2 & log & 644 & 2 & 0 & 67/69 \\
github-copilot & 2 & log/state & 600/644 & 0 & n/a & 236/248 \\
grok-build & 2 & log & 644 & 0 & 0 & 35/54 \\
opencode & 0 & -- & -- & 0 & n/a & 26/26 \\
qwen-code & 1 & log & 644 & 1 & 0 & 10/13 \\
\bottomrule
\end{tabular}
\end{adjustbox}
\end{table}

\subsection{E4: Distribution supply chain (RQ2, D13)}\label{sec:eval-supply}
We inspected the published package of \eFourHarnesses{} harnesses (the Tier-A set plus Claude Code, Grok Build and Copilot CLI) from registry metadata and lockfiles only, without running install scripts\evD{} (per-package table: \texttt{paper/tables/e4\_supply.md}). \eFourHarnessesWithInstallHooks{} packages declare install-time hooks. \eFourWithProvenance{} of \eFourProvenanceKnown{} with determinable status publish registry provenance attestations. \eFourHarnessesWithCurlInstaller{} offer a \texttt{curl | sh} installer, and in \eFourHarnessesWithUnsignedCurlInstaller{} of them the installer verifies no signature (some verify a checksum; Claude Code documents a signed release manifest that its installer does not itself verify\evD). Auto-update is documented as on by default in \eFourAutoUpdateDefaultOn{} (cells marked ``unverified'' in the table are undetermined, not off). The median lockfile size is \eFourMedianTransitive{} transitive packages, and OSV lists \eFourOsvUniqueVulnsTotal{} unique advisories against locked dependency versions, concentrated in a few packages. These are advisories against pinned versions, not demonstrated exploitability. Auto-update is a trade-off: it delivers fixes quickly but widens the blast radius of a compromised release channel, as the Amazon Q extension incident showed\evD~\cite{amazonq2025}.

